% §3 Problem definition.  One continuous argument: the STDR recursion, the Theta(m^2) cost of its
% top-level call, the Bernoulli(p) observation model that answers that cost, and the entry-wise
% criterion eq:entrywise that every bound in the paper is measured against.  No subsections --
% the section is one argument and headings would restart it.
\section{Problem definition}\label{sec:problem}

Because the Fiedler split of the full similarity matrix is always a clan pair
(\Cref{thm:stdr-split}), tree reconstruction becomes a divide-and-conquer recursion. Spectral
Top-Down Reconstruction (STDR) \cite{aizenbud2023spectral} forms the similarity matrix over the
current node set, splits it by the sign pattern of the Fiedler vector, solves each side
independently by any accurate reconstruction method once the subproblem is small enough, and merges
the two subtrees along the recovered edge (\Cref{alg:stdr}).

\begin{algorithm}[t]
\caption{Spectral Top-Down Reconstruction (STDR) \cite{aizenbud2023spectral}}
\label{alg:stdr}
\begin{algorithmic}[1]
\Procedure{STDR}{$\mathcal{X}$}
  \If{$|\mathcal{X}| \le m_0$}
    \State \Return $\textsc{BaseReconstruct}(\mathcal{X})$
           \Comment{any accurate method; superlinear in $|\mathcal{X}|$}
  \EndIf
  \State $S \gets \bigl[\,S(x_i,x_j)\,\bigr]_{x_i,x_j \in \mathcal{X}}$
         \Comment{$\Theta(|\mathcal{X}|^2)$ pairwise estimates}
  \State $L \gets D - S$ \Comment{$D_{ii}=\sum_j S_{ij}$, \Cref{def:fiedler}}
  \State $v \gets$ eigenvector of $L$ at $\lambda_2$ \Comment{the Fiedler vector}
  \State $C_1 \gets \{x_i : v_i \ge 0\}$, \quad $C_2 \gets \{x_j : v_j < 0\}$
         \Comment{adjacent clans, \Cref{thm:stdr-split}}
  \State $\mathcal{T}_1 \gets \textsc{STDR}(C_1)$, \quad
         $\mathcal{T}_2 \gets \textsc{STDR}(C_2)$
  \State \Return $\textsc{Merge}(\mathcal{T}_1, \mathcal{T}_2)$
         \Comment{join the two subtrees at the recovered edge}
\EndProcedure
\end{algorithmic}
\end{algorithm}

The recursion is attractive because the expensive part of tree reconstruction is superlinear in the
number of terminal nodes, so cutting the problem in half repeatedly reduces the total work sharply,
while \Cref{thm:stdr-split} guarantees that each cut is a genuine split and the merge is therefore
well-posed \cite{aizenbud2023spectral}.

The top-level call of \Cref{alg:stdr} is the one that governs both the error and the cost. It is the
only call that sees all $m$ terminal nodes, and no parent call corrects it: a wrong split at the
root cannot be repaired further down. It is also where the similarity matrix is largest.

That is the bottleneck, and it sits on line 5. An entry of $S$ is not a table lookup but a statistic
of a pair of sequences: the similarity of \Cref{def:similarity} is estimated from the joint
character counts of terminal nodes $i$ and $j$ over a length-$\ell$ alignment. Forming the full
matrix at the top level therefore costs
\begin{equation}\label{eq:cost}
   \binom{m}{2} = \Theta(m^2) \ \text{pairwise estimates},
   \qquad \Theta(m^2\ell) \ \text{character comparisons},
\end{equation}
which for $m \sim 10^5$ terminal nodes is on the order of $10^{10}$ pairwise estimates before the
recursion has taken a single step. The quadratic term is exactly the cost that divide-and-conquer
was introduced to avoid, and it is paid in full at the root.

This raises the question the paper answers: can the top-down split be computed at almost linear
cost, with a guarantee that it is still the same split the full matrix would have returned?

Our answer is to compute only a random subset of the pairs. Following the matrix-completion
paradigm \cite{candes2009power, candes2008exact}, we observe entries through a random mask and
reweight: let $\Omega\in\{0,1\}^{m\times m}$ be a symmetric mask whose upper-triangular entries are
independent $\mathrm{Bernoulli}(p)$, $p\in(0,1]$, and form the empirical similarity by
inverse-probability weighting \cite{candes2009matrix},
\begin{equation}\label{eq:sampling_estimator}
   \hat S_{ij}=\begin{cases}\tfrac1p\,\Omega_{ij}\,S_{ij}, & i\ne j,\\ 1, & i=j.\end{cases}
\end{equation}
Only the $\Theta(p m^2)$ pairs in the mask are ever estimated, so $p$ is the fraction of the cost in
\eqref{eq:cost} that is actually paid. Under the assumptions of \Cref{sec:assumptions}, a rate
$p=\Theta(\log m/m)$ suffices to reproduce the full-data split (\Cref{thm:main-sim}), which brings
the top-level call down from $\Theta(m^2)$ to $\Theta(m\log m)$ pairwise estimates.

This estimator is unbiased, $\E[\hat S]=S$, so $\EL=\hat L-L$ is mean zero, where $\hat L=\hat
D-\hat S$. The resulting error is missingness rather than additive noise: most pairs are left at
zero and the rest rescaled by $1/p$, so it is heavy, structured, and its per-node variance grows as
$p\to0$ (\Cref{prop:C-var}).

\medskip
\noindent\textbf{Problem.} Given the sparsely observed estimate $\hat S$ of
\eqref{eq:sampling_estimator}, recover the Fiedler split $\Pisp{S}$ of the full matrix from the sign
pattern of the Fiedler vector $\vvh$ of $\hat L$. We ask for the minimal sampling rate $p$ at which
$\Pisp{\hat S}=\Pisp{S}$ \hp. Since $\E[\hat S]=S$, the question reduces to an
eigenvector-perturbation problem over $\EL=\hat L-L$, which the remainder of the paper analyzes.
The full-data split $\Pisp{S}$ is already a clan pair of $\mathcal{T}$ (\Cref{thm:stdr-split}), so a
rate at which $\Pisp{\hat S}=\Pisp{S}$ is a rate at which the sub-sampled operator returns a genuine
split of the tree.

\medskip
Recovery is, at its core, an \emph{entry-wise} event. Two vectors induce the same
sign partition exactly when they agree in sign at every node, so the top-down
partition succeeds precisely when the empirical Fiedler vector agrees in sign with
its population counterpart coordinate by coordinate, that is, when
\begin{equation}\label{eq:entrywise}
   \norm{\vvh-\vv}_\infty<\min_i|\vv_i|.
\end{equation}
Condition \eqref{eq:entrywise} is the formal recovery event that every bound in the paper is
measured against. Its right-hand side is not computable without a model for $S$;
\Cref{sec:assumptions} supplies one, under which $\min_i|\vv_i|$ has a closed form
(\Cref{lem:A-decomp}).

The scope of the analysis is the population matrix. The empirical estimator of an individual entry
concentrates around $S_{ij}$ at rate $\mathcal O(1/\sqrt\ell)$ \cite{aizenbud2023spectral}, and the
structural conditions of \Cref{sec:assumptions} are stated as properties of the population $S$,
whose observable proxy is that estimator.
